\documentclass[aspectratio=169,11pt]{beamer}

\usetheme{Madrid}
\usecolortheme{default}
\usefonttheme{professionalfonts}
\setbeamertemplate{navigation symbols}{}
\setbeamertemplate{footline}[frame number]

\usepackage{amsmath, amssymb, booktabs}

\title{Spatial Mismatch, Segregation, Neighborhood Effects}
\subtitle{And Job Access}
\author{Special Topic 4: Urban and Labor -- Week 3}
\date{}

\begin{document}

\begin{frame}
  \titlepage
\end{frame}

\begin{frame}{Central question}
\begin{itemize}
    \item How do residential segregation and neighborhood exposure shape access to jobs, networks, and long-run labor-market opportunity?
    \item This is a labor question because place changes search, hiring, commuting, networks, skills, earnings, employment, and mobility.
    \item Week 3 discipline: separate mechanisms before interpreting neighborhood or segregation gaps.
    \item The goal is unequal access, not aggregate access.
\end{itemize}
\end{frame}

\begin{frame}{From aggregate access to unequal access}
\begin{itemize}
    \item Week 1: local labor markets are systems of jobs, residences, commuting costs, rents, amenities, and outside options.
    \item Week 2: housing and moving costs can block access to high-opportunity places.
    \item Week 3: workers in the same metro area can face different feasible job sets.
    \item Strong metro labor demand can coexist with weak access for segregated or transit-constrained neighborhoods.
\end{itemize}
\end{frame}

\begin{frame}{Job-access object}
\[
A_i = \sum_j v_j \exp\{-\kappa \tau_{ij}\}
\]
\small
\begin{tabular}{p{0.20\linewidth} p{0.64\linewidth}}
\toprule
Term & Interpretation \\
\midrule
$A_i$ & effective job access for worker or residence $i$ \\
$v_j$ & job opportunities, vacancies, or job quality at location $j$ \\
$\tau_{ij}$ & generalized travel cost: time, money, reliability, transfers, safety, schedules \\
$\kappa$ & rate at which access decays with travel cost \\
\bottomrule
\end{tabular}
\end{frame}

\begin{frame}{Mechanism taxonomy}
\small
\begin{tabular}{p{0.32\linewidth} p{0.46\linewidth}}
\toprule
Mechanism & Main labor object \\
\midrule
Distance to jobs & reachable vacancies and commute feasibility \\
Transit and schedules & generalized travel cost and feasible shifts \\
Address stigma & employer callbacks, interviews, offers \\
Residential networks & referrals, information, coworker links \\
Peers and social capital & search behavior, aspirations, expectations \\
Schools and institutions & human capital and long-run opportunity \\
Safety and exposure & work timing, search radius, retention, welfare \\
\bottomrule
\end{tabular}
\end{frame}

\begin{frame}{Mechanism decomposition}
\[
Y_i = \beta_A A_i + \beta_N N_i + \beta_P P_i + \beta_S S_i + u_i
\]
\begin{itemize}
    \item $A_i$: job access and commuting feasibility.
    \item $N_i$: referral networks and neighbor-coworker links.
    \item $P_i$: peers, social capital, safety, local expectations.
    \item $S_i$: schools, public goods, and local institutions.
    \item The warning: one neighborhood coefficient can bundle several mechanisms.
\end{itemize}
\end{frame}

\begin{frame}{Spatial mismatch, networks, neighborhood effects}
\small
\begin{tabular}{p{0.25\linewidth} p{0.30\linewidth} p{0.30\linewidth}}
\toprule
Channel & Counterfactual & Natural design \\
\midrule
Spatial mismatch & same worker, lower travel cost to jobs & travel-time, transit, job-location shocks \\
Networks & same worker, different referral environment & matched worker-firm and neighbor links \\
Address stigma & same applicant, different address signal & audit or correspondence experiment \\
Neighborhood exposure & same child, different duration or timing of exposure & mover or policy-induced relocation design \\
\bottomrule
\end{tabular}
\end{frame}

\begin{frame}{Job suburbanization and unequal access}
\begin{itemize}
    \item Jobs can move even when residential segregation persists.
    \item A suburban job center may raise total metro employment but lower effective access for transit-dependent central-city residents.
    \item The labor margin may be employment, commute burden, job quality, or search scope.
    \item Proximity is not access: employers, referrals, schedules, and transport all mediate the connection.
\end{itemize}
\end{frame}

\begin{frame}{Neighborhood exposure and long-run outcomes}
\[
Y_i = \alpha + \theta \, \text{Exposure}_{ig} + \lambda_g + \varepsilon_i
\]
\begin{itemize}
    \item Exposure designs ask whether duration or age at exposure changes adult outcomes.
    \item Earlier moves can reveal causal place effects when timing is plausibly comparable.
    \item Labor outcomes include later earnings, employment, college, occupational mobility, and search scope.
    \item Mechanisms may run through schools, peers, safety, institutions, networks, or expectations.
\end{itemize}
\end{frame}

\begin{frame}{Mover designs and policy shocks}
\small
\begin{tabular}{p{0.28\linewidth} p{0.34\linewidth} p{0.25\linewidth}}
\toprule
Design & Best for & Main threat \\
\midrule
Mover / age-at-move & long-run exposure timing & endogenous moves and disruption \\
Cohort exposure & cumulative duration and age profiles & omitted family trends \\
MTO or vouchers & neighborhood opportunity change & treatment bundles \\
Public housing demolition & policy-induced relocation & displacement channels \\
Transit or job-location shock & current access change & equilibrium resorting \\
\bottomrule
\end{tabular}
\end{frame}

\begin{frame}{School-boundary, audit, matched-firm designs}
\small
\begin{tabular}{p{0.30\linewidth} p{0.42\linewidth} p{0.16\linewidth}}
\toprule
Design & Best use & Key limit \\
\midrule
School boundary or admission & separate schools from neighborhoods & capitalization \\
Audit or correspondence & address stigma, distance, employer beliefs & callback margin \\
Matched employer-employee & referrals, coworker geography, workplace segregation & data access \\
Neighbor-coworker overlap & residential labor-market networks & sorting \\
\bottomrule
\end{tabular}
\end{frame}

\begin{frame}{What good empirical design does}
\begin{itemize}
    \item Maps the treatment to one main mechanism.
    \item Defines geography using commuting, schools, institutions, or employer markets.
    \item Measures actual access, exposure, or employer response.
    \item Distinguishes residence, workplace, school, and neighborhood channels.
    \item Addresses residential sorting and timing.
    \item Discusses resorting, displacement, and capitalization when relevant.
    \item States which labor margin moves.
\end{itemize}
\end{frame}

\begin{frame}{Data sources}
\small
\begin{tabular}{p{0.31\linewidth} p{0.50\linewidth}}
\toprule
Source & Typical use \\
\midrule
LEHD / LODES / QWI & residence-workplace flows and job geography \\
Census / ACS / Decennial & segregation, commuting, earnings, sorting \\
Opportunity Atlas / tax-linked data & childhood exposure and adult outcomes \\
School boundaries / assignment & school versus neighborhood channels \\
GTFS / travel-time matrices & generalized job access \\
Audit data & employer response to address or commute signals \\
Matched employer-employee data & referrals and workplace segregation \\
Housing or relocation admin data & policy-induced moves and exposure \\
\bottomrule
\end{tabular}
\end{frame}

\begin{frame}{Week 3 lab}
\begin{itemize}
    \item Reproduce: synthetic age-at-move exposure logic anchored in MTO.
    \item Diagnose: classify mechanisms and identify the counterfactual.
    \item Transfer: apply the same discipline to job suburbanization.
    \item Goal: separate access from employment, neighborhood effects from sorting, schools from neighborhoods, and policy shocks from mover designs.
\end{itemize}
\end{frame}

\begin{frame}{Bridge to Week 4}
\begin{itemize}
    \item Week 3 treats safety as one mechanism inside unequal access.
    \item Week 4 studies safety and environmental exposure directly as constraints on work.
    \item The same design discipline carries forward: define the labor margin, mechanism, geography, and counterfactual.
\end{itemize}
\end{frame}

\begin{frame}{Takeaways}
\begin{itemize}
    \item Segregation matters for labor when it changes access, networks, stigma, schools, exposure, safety, or institutions.
    \item Proximity to jobs is not the same as access to jobs.
    \item Neighborhood effects require designs that confront sorting and channel bundling.
    \item The best papers identify a mechanism and a labor margin, not a generic place effect.
\end{itemize}
\end{frame}

\end{document}
